You will build the classic first rocket --- a nose cone, a body tube, and three
fins --- verify it, fly it on a B6, and read every line of the flight report.

\begin{note}[Setup]
Work in an empty directory containing a copy of the small Estes-class motor
database that ships with QtRocket:
\begin{minted}{text}
mkdir flight1 && cd flight1
cp <qtrocket>/tests/data/motors/estes_small.qmd .
qtrocket-cli
\end{minted}
The complete session below is also included as
\srcf{scripts/tutorial1.cli}, runnable with
\code{qtrocket-cli tutorial1.cli}.
\end{note}

\subsection*{Load motors, find one}

\begin{console}
\muted{# QtRocket CLI - type 'help' for commands, 'quit' to exit.}
\uin{loaddb estes_small.qmd}
\ok loaddb: 13 new motors from estes_small.qmd (13 total)
\uin{listmotors B}
\ok listmotors: 2 shown (filter="B")
B4  avg=4.17N  Itot=5Ns
B6  avg=5.03N  Itot=4.33Ns
\end{console}


Two commands, two habits. \repl{loaddb} fills the session motor database from a
\code{.qmd} file and answers with how many motors arrived --- every command
answers \ok{} or \err{}, and that is the entire protocol. \repl{listmotors}
takes an optional substring filter (case-sensitive); \code{B} narrows the
thirteen bundled motors to the B-class. The two numbers per row --- average
thrust and total impulse --- are the two numbers that matter when choosing a
motor.

\subsection*{Build the airframe}

\begin{console}
\uin{newdesign NoseCone name=nose baseRadius=0.0098 length=0.10 density=700}
\ok newdesign: root NoseCone id=2
\uin{addpart nose BodyTube name=body innerRadius=0.0091 outerRadius=0.0098 length=0.25 density=700}
\ok addpart: BodyTube id=3 under 2
\uin{addpart body FinSet name=fins finCount=3 rootChord=0.04 tipChord=0.02 span=0.03 sweep=0.02 thickness=0.003 bodyRadius=0.0098 density=700 seat=OnSurface parentStation=0 childStation=0}
\ok addpart: FinSet id=4 under 3
\uin{checkdesign}
\ok checkdesign: 3 parts, contiguous, no overlaps
\uin{listparts}
\ok listparts:
  [2] NoseCone "nose"  m=0.0070401 kg
    [3] BodyTube "body"  m=0.00727357 kg
      [4] FinSet "fins"  m=0.00567 kg
  -- composite: mass=0.0199837 kg, cg_z=-0.200687 m (t=0)
\end{console}


Read the \repl{addpart} lines closely; they carry the whole design grammar:

\begin{itemize}
\item \repl{newdesign} starts the tree with a root part --- here a solid
      conical nose, $0.0098\m$ base radius, $0.10\m$ long, at
      $700$~kg/m$^3$ (a plausible balsa-and-paper blend).
\item The body tube names its parent: \code{addpart \textbf{nose} BodyTube
      ...}. Parents are addressed by the \code{name=} you gave them --- names,
      not ids, are the scriptable handle. No placement keys are needed: the
      default seat (\seat{Abut}) hangs the tube's fore rim flush under the
      nose's aft plane.
\item The fins use \code{seat=OnSurface} with
      \code{parentStation=0 childStation=0}: sit on the body's outer wall,
      aft edges aligned --- fins flush with the bottom of the tube.
\end{itemize}

\repl{checkdesign} then renders the verdict that matters:
\code{3 parts, contiguous, no overlaps}. \repl{listparts} shows the tree with
each part's computed mass and the composite mass and center of gravity ---
derived from geometry and density, not typed in.

\subsection*{Fly it}

\begin{console}
\uin{setmotor B6}
\ok setmotor: B6
\uin{setdrag 0.75}
\ok setdrag: 0.75
\uin{setarea 0.000302}
\ok setarea: 0.000302 m^2
\uin{status}
\ok status:
  motor      = B6
  dry_mass   = 1 kg
  drag_coeff = 0.75
  ref_area   = 0.000302 m^2
  velocity   = 0 m/s
  angle      = 0 deg (from vertical)
  atmosphere = Constant Atmosphere
  gravity    = Constant Gravity
  integrator = Runge-Kutta 4th Order
  database   = 13 motors
\uin{launch}
\ok launch: 1457 steps, t_final=14.560s
  apogee    = 241.405 m @ t=6.330s
  max_speed = 97.562 m/s @ t=0.820s
  downrange = 0.000 m
  landing   = t=14.560s, pos=(0.000, 0.000, -0.314)
CSV /tmp/flight1/qtrocket_run.csv
\end{console}


\repl{setmotor} attaches the B6 by its common name. \repl{setdrag} and
\repl{setarea} set the drag coefficient and aerodynamic reference area ---
$0.75$ and the body tube's frontal disc ($\pi r^2 \approx 3.02\times10^{-4}$~m$^2$)
are honest starting values for a small sport model. \repl{status} plays back
the full session configuration before committing to flight. One row deserves a
squint: \code{dry\_mass = 1 kg} is only the mirror of the last \repl{setmass}
--- an initial placeholder here, never applied, since this session hasn't
called \repl{setmass} at all. Your rocket flies on the part tree's computed
$\approx 20$\,g composite plus the motor, as \repl{listparts} showed. (On a
design, \repl{setmass} immediately overwrites the \emph{root part's} own mass
--- rarely what you want on a multi-part rocket; trust the tree.)

The \repl{launch} report is the payoff, one line per question you would ask at
the field: how high (\code{apogee}), how fast (\code{max\_speed}, reached just
after burnout), how far downwind (\code{downrange} --- zero, this was a
straight-up flight), and where and when it came down. The \code{CSV} line
names the full trajectory file --- every timestep, every state --- written for
plotting or analysis (\cref{sec:tut5} uses it).

\subsection*{Keep it}

\begin{console}
\uin{savedesign first_rocket.qrd}
\ok savedesign: first_rocket.qrd
\uin{quit}
\ok bye
\end{console}


\repl{savedesign} writes the rocket --- tree, placement, drag config, and the
motor by name --- to a \code{.qrd} file. Open it in the visualizer to see what
you built (\cref{fig:first-rocket}):

\begin{minted}{text}
qtrocket-visualizer first_rocket.qrd
\end{minted}

\screenshot{ug-first-rocket.png}{Your first rocket in
\code{qtrocket-visualizer}: conical nose, 18\,mm-class body, three swept fins
flush with the aft rim. Orbit with the left mouse button, zoom with the
wheel.}{fig:first-rocket}{0.82}

\begin{tip}[What you learned]
\repl{loaddb} $\rightarrow$ \repl{listmotors} $\rightarrow$ \repl{newdesign}
$\rightarrow$ \repl{addpart} (names as handles, seats as intent) $\rightarrow$
\repl{checkdesign} $\rightarrow$ \repl{setmotor} $\rightarrow$ \repl{launch}
$\rightarrow$ \repl{savedesign}. That eight-step spine is every QtRocket
session; the rest of this guide is variations.
\end{tip}
